% !TeX root = ../main.tex
\chapter{Phân tích cây (Tree Analysis)}
\label{chap:tree-analysis}

\section{Cấu trúc cây}
Cây baseline (cây gốc) khá phức tạp với độ sâu (depth) lên đến 7 tầng, bao gồm tổng cộng 19 leaf và 37 node. Sự split sâu và rườm rà này thường là nguyên nhân khiến cây dễ bị phụ thuộc quá mức vào dữ liệu huấn luyện.

\section{Root/Key Splits}
Đặc trưng (feature) được chọn ở node gốc (root) và những lần split đầu tiên đóng vai trò quan trọng nhất trong việc phân tách ranh giới giữa khối u ác tính và lành tính. Cụ thể, node gốc sử dụng ngưỡng phân chia \texttt{worst radius $\le$ 16.795} (trung bình của 3 giá trị bán kính lớn nhất của các nhân tế bào). Split đầu tiên này giúp phân tách phần lớn các mẫu lành tính (Benign) có kích thước nhân nhỏ sang một nhánh riêng, làm giảm mạnh độ pha tạp (Impurity / Gini). Bên cạnh đó, đặc trưng \texttt{worst concave points} (số lượng và mức độ nghiêm trọng của các điểm lõm trên đường viền nhân tế bào) cũng xuất hiện ở những lần split đầu tiên. Về mặt y khoa, nhân tế bào ung thư thường có hình dạng méo mó, đường viền gồ ghề và nhiều điểm lõm bất thường so với tế bào lành tính tròn đều, do đó đây là những đặc trưng mang tính quyết định cao.

\section{Dấu hiệu Overfitting / Underfitting}
Bằng chứng rõ rệt về hiện tượng overfitting (học vẹt) thể hiện qua việc mô hình Baseline đạt độ chính xác tuyệt đối 100\% trên tập huấn luyện (Train Accuracy = 1.0), nhưng độ chính xác trên tập kiểm tra (Test Accuracy) lại tụt xuống chỉ còn 91.23\%. Sự chênh lệch lớn này, kết hợp với cấu trúc phức tạp 19 leaf và độ sâu 7 tầng (các node leaf kích thước siêu nhỏ chỉ chứa 1 hoặc 2 mẫu xuất hiện dày đặc ở các tầng sâu từ tầng 5 đến 7), chứng minh cây đang cố gắng phân tách triệt để các điểm ngoại lai (outliers/nhiễu) của tập train, làm giảm đi khả năng tổng quát hóa. Trực quan về cây baseline được thể hiện qua Hình~\ref{fig:baseline-readable}.
